% Lernkarten -- Analysis (Modul 61211), Lektion 2
% Quelle: eigene OneNote-Notizen (Fernuni Notizbuch > 61211 Analysis)
% Format: \lernkarte{Vorderseite (roter Titel)}{Rueckseite (weisser Inhalt)}

% --- OneNote-Seite 2.1 "Rückblick Stetigkeit" ---------------------------
\lernabschnitt{2.1 Rückblick Stetigkeit}

\lernkarte[61211-2-01]{Definition Umkehrfunktion}{%
  $f : A \to B$. \\
  $f$ ist injektiv, dann $f^{-1} : f(A) \to B$.
}

\lernkarte[61211-2-02]{Definition Stetigkeit}{%
  \[ \forall V \text{ Umgebung von } f(a) \ \exists U \text{ Umgebung von } a:
     \ f(U \cap M) \subseteq V \]
}

\lernkarte[61211-2-03]{Restriktion / lokale Eigenschaft}{%
  $f|_{U\cap M}$ ist stetig. \\
  $f|_{U\cap M}$ stetig in $a \;\Rightarrow\; f$ stetig in $a$.
}

\lernkarte[61211-2-04]{Stetigkeitskriterien}{%
  \begin{itemize}[topsep=2pt,itemsep=3pt,leftmargin=*]
    \item \textbf{$\varepsilon$-$\delta$-Kriterium:}
      $\forall\varepsilon>0 \ \exists\delta>0 \ \forall x\in M:
       |x-a|<\delta \Rightarrow |f(x)-f(a)|<\varepsilon$
    \item \textbf{Folgenkriterium:} $\forall (x_k)$ mit
      $\lim_{k\to\infty} x_k = a: \ \lim_{k\to\infty} f(x_k) = f(a)$
    \item \textbf{Kettenregel:} $f$ stetig in $a \wedge g$ stetig in $f(a)
      \Rightarrow (g\circ f)$ ist stetig in $a$
  \end{itemize}
}

\lernkarte[61211-2-05]{lokale Beschränktheit}{%
  $f|_{U\cap M}$ (Restriktion) ist beschränkt.
}

% PRÜFEN: Notiz schwer lesbar ("f(a) < f(b)"); sinngemäß Trennung zweier
% Funktionen -- bitte gegen OneNote prüfen.
\lernkarte[61211-2-06]{lokale Trennung}{%
  \[ f(a) < g(a) \;\Rightarrow\; \forall x \in U\cap M: \ f(x) < g(x) \]
}

\lernkarte[61211-2-07]{Zwischenwertsatz}{%
  \[ I = [y_1, y_2], \quad \forall c \in (y_1, y_2) \ \exists \xi: \ f(\xi) = c \]
}

\lernkarte[61211-2-08]{Eigenschaften strenge Monotonie}{%
  \begin{itemize}[topsep=2pt,itemsep=1pt,leftmargin=*]
    \item injektiv $\Leftrightarrow$ streng monoton
    \item streng monoton $\Rightarrow$ stetige, streng monotone Umkehrfunktion
  \end{itemize}
}

\lernkarte[61211-2-09]{Definition Häufungspunkt}{%
  Wenn ein $x \neq a$ in jeder Umgebung von $a$ liegt. \\
  $\forall x$ ist Häufungspunkt in $M \;\Rightarrow\; M$ abgeschlossen.
}

\lernkarte[61211-2-10]{Satz von Bolzano-Weierstraß (Häufungspunkte)}{%
  \[ M \text{ beschränkt} \;\Rightarrow\; \exists \text{ Häufungspunkt} \]
}

\lernkarte[61211-2-11]{Kompakte Teilmenge von $\mathbb{R}$}{%
  \[ \text{kompakt} \;\Leftrightarrow\; \text{abgeschlossen} \wedge \text{beschränkt} \]
}

\lernkarte[61211-2-12]{Stetiges Bild einer kompakten Menge}{%
  \[ M \text{ kompakt} \;\Rightarrow\; \text{Bild von } M \text{ ist kompakt} \]
}

\lernkarte[61211-2-13]{Satz von Weierstraß}{%
  $M$ kompakt $\Rightarrow \exists$ Maximum und Minimum; gleichmäßig stetig.
}

\lernkarte[61211-2-14]{gleichmäßige Stetigkeit}{%
  \[ \forall\varepsilon>0 \ \exists\delta>0 \ \forall x,y\in M:
     \ |x-y|<\delta \Rightarrow |f(x)-f(y)|<\varepsilon \]
}

% --- OneNote-Seite 2.3 "Stetigkeit R^n" ---------------------------------
\lernabschnitt{2.3 Stetigkeit $\mathbb{R}^n$}

\lernkarte[61211-2-15]{Definition Stetigkeit (in $\mathbb{R}^n$)}{%
  \[ \forall V \text{ Umgebung von } f(a) \ \exists U:
     \ f(U \cap M) \subseteq V \]
}

\lernkarte[61211-2-16]{Stetigkeitskriterien (in $\mathbb{R}^n$)}{%
  \begin{itemize}[topsep=2pt,itemsep=3pt,leftmargin=*]
    \item \textbf{$\varepsilon$-$\delta$-Kriterium:}
      $\forall\varepsilon>0 \ \exists\delta>0 \ \forall x\in M:
       \|x-a\|_x<\delta \Rightarrow \|f(x)-f(a)\|_y<\varepsilon$
    \item \textbf{Folgenkriterium:} $x_k \to a \Rightarrow f(x_k) \to f(a)$
    \item \textbf{Kettenregel:} $f$ stetig in $a \wedge g$ stetig in $f(a)
      \Rightarrow g\circ f$ stetig in $a$
    \item \textbf{Dehnungsbeschränkt:} $\forall x,x':
      \|f(x)-f(x')\| \le L\,\|x-x'\|$
  \end{itemize}
}

\lernkarte[61211-2-17]{Restriktion / lokale Eigenschaft (in $\mathbb{R}^n$)}{%
  \begin{itemize}[topsep=2pt,itemsep=1pt,leftmargin=*]
    \item $f$ in $a$ stetig $\Rightarrow$ Restriktion $f|_U$ stetig in $a$
    \item $f|_U$ stetig mit Umgebung $U \Rightarrow f$ ist stetig in $a$
  \end{itemize}
}

\lernkarte[61211-2-18]{lokale Beschränktheit (in $\mathbb{R}^n$)}{%
  $f$ stetig in $a \;\Rightarrow\; f|_{U\cap M}$ ist beschränkt.
}

\lernkarte[61211-2-19]{lokale Trennung (in $\mathbb{R}^n$)}{%
  \[ f(a) \neq g(a) \;\Rightarrow\; f(x) \neq g(x) \quad \forall x \in U(a) \]
}

% --- OneNote-Seite 2.4 "Zusammenhängende Mengen" ------------------------
\lernabschnitt{2.4 Zusammenhängende Mengen}

\lernkarte[61211-2-20]{Definition offene Menge}{%
  \[ M \text{ ist offen} := \forall x \in M: \ M \text{ ist Umgebung von } x \]
}

\lernkarte[61211-2-21]{Definition abgeschlossene Menge}{%
  Jeder Häufungspunkt liegt in $M \;\Rightarrow\; M$ ist abgeschlossen.
}

\lernkarte[61211-2-22]{Topologische Stetigkeit}{%
  \[ \forall S \text{ offen} \ \exists Q \text{ offen}: \ f^{-1}(S) = Q \cap M \]
}

\lernkarte[61211-2-23]{Definition zusammenhängende Mengen}{%
  Keine Trennung in zwei Mengen durch zwei getrennte offene Mengen.
}

\lernkarte[61211-2-24]{Zusammenhängende Mengen in $\mathbb{R}$}{%
  \[ \text{zusammenhängend} \;\Leftrightarrow\; \text{Intervall} \]
}

\lernkarte[61211-2-25]{Stetiges Bild zusammenhängender Menge}{%
  \[ M \text{ zusammenhängend} \;\Rightarrow\; f(M) \text{ ist zusammenhängend} \]
}

% --- OneNote-Seite 2.5 "Kompakte Menge" ---------------------------------
\lernabschnitt{2.5 Kompakte Menge}

\lernkarte[61211-2-26]{Definition Kompaktheit}{%
  Jede offene Überdeckung besitzt eine (endliche) Teilüberdeckung.
}

\lernkarte[61211-2-27]{Folgenkompaktheit}{%
  Jede Folge besitzt eine konvergente Teilfolge.
}

\lernkarte[61211-2-28]{Folgerungen Kompaktheit}{%
  \begin{itemize}[topsep=2pt,itemsep=1pt,leftmargin=*]
    \item kompakt $\Rightarrow$ abgeschlossen $\wedge$ beschränkt
    \item stetig: nimmt Maximum und Minimum an
  \end{itemize}
}

\lernkarte[61211-2-29]{Satz von Heine-Borel}{%
  \[ \text{in } \mathbb{R}^n: \quad \text{kompakt} \;\Leftrightarrow\;
     \text{abgeschlossen} \wedge \text{beschränkt} \]
}

\lernkarte[61211-2-30]{Stetiges Bild kompakter Menge}{%
  \[ f \text{ stetig} \wedge M \text{ kompakt} \;\Rightarrow\;
     f(M) \text{ ist kompakt} \]
}

\lernkarte[61211-2-31]{Definition \& Satz gleichmäßige Stetigkeit}{%
  \[ \forall\varepsilon>0 \ \exists\delta>0 \ \forall x,y\in M:
     \ \|x-y\|<\delta \Rightarrow \|f(x)-f(y)\|<\varepsilon \]
  $M$ kompakt und $f$ stetig $\Rightarrow f$ gleichmäßig stetig.
}

% --- OneNote-Seite 2.6 "Konvergenz Funktionenfolgen" --------------------
\lernabschnitt{2.6 Konvergenz Funktionenfolgen}

\lernkarte[61211-2-32]{Definition punktweise Konvergenz}{%
  $f_k(x) \to f(x)$ mit $k \to \infty$ für jedes $x$.
}

\lernkarte[61211-2-33]{Definition gleichmäßige Konvergenz}{%
  \[ \sup\{\, \|f_k(x) - f(x)\| \;\mid\; x \in M \,\} \;\longrightarrow\; 0 \]
}

\lernkarte[61211-2-34]{$\varepsilon$-$n_0$-Definition der Konvergenz}{%
  \begin{itemize}[topsep=2pt,itemsep=3pt,leftmargin=*]
    \item \textbf{punktweise:}
      $\forall x\in M \ \forall\varepsilon>0 \ \exists n_0\in\mathbb{N} \ \forall k\in\mathbb{N}:
       k \ge n_0 \Rightarrow \|f_k(x)-f(x)\|<\varepsilon$
    \item \textbf{gleichmäßig:}
      $\forall\varepsilon>0 \ \exists n_0\in\mathbb{N} \ \forall x\in M \ \forall k\in\mathbb{N}:
       k \ge n_0 \Rightarrow \|f_k(x)-f(x)\|<\varepsilon$
  \end{itemize}
}

% PRÜFEN: Notiz schreibt "gleichmäßig stetig"; der Satz von Weierstraß
% braucht gleichmäßige KONVERGENZ (so hier gesetzt) -- bitte prüfen.
\lernkarte[61211-2-35]{Satz von Weierstraß (Funktionenfolgen)}{%
  $f_k$ stetig in $a \;\wedge\; (f_k)$ gleichmäßig konvergent
  $\;\Rightarrow\; f$ ist stetig in $a$.
}

\lernkarte[61211-2-36]{Definition konvergente Funktionenreihe}{%
  Funktionenreihe: $\displaystyle\sum_{k=1}^{\infty} f_k$. \\
  $(S_k)$ ist punktweise bzw.\ gleichmäßig konvergent.
}

\lernkarte[61211-2-37]{Stetigkeit der Summenfunktion}{%
  \[ \sum_{k=1}^{\infty} f_k \ \xrightarrow{\text{gleichm.}}\ S
     \ \wedge\ f_k \text{ stetig in } a \;\Rightarrow\; S \text{ stetig in } a \]
}

\lernkarte[61211-2-38]{Majorantenkriterium}{%
  $\forall k \ \exists a_k$ (eine Schranke): $\ \sup\{\|f_k(x)\| \mid x\in M\} \le a_k$
  \[ \wedge \ \sum_{k=1}^{\infty} a_k \text{ konvergent}
     \;\Rightarrow\; \sum_{k=1}^{\infty} f_k \text{ konvergiert gleichmäßig} \]
}
